\section{Results}

The current results are validation-stage results rather than final held-out
generation results. They show that the implemented training pipeline is able to
train xLSTM language models on REMI-style music tokens, compare checkpoints, and
identify stronger configurations for follow-up evaluation. The results should
therefore be interpreted as model-selection evidence, not as a final statement
about musical quality.

\subsection{Validation Results}

Table~\ref{tab:result-checkpoints} summarizes representative completed runs from
the local \verb|summary.json| files. The strongest validation-loss checkpoint is
the base xLSTM with tokenizer 14 and batch size 16. The batch-size-8 tokenizer-14
run is almost tied and keeps the same practical setup used by the broader sweep.
The medium model does not improve the validation-selected result under its
current 15-epoch restart schedule, even though it has more than twice as many
parameters.

\begin{table}[htbp]
	\centering
	\scriptsize
	\resizebox{\textwidth}{!}{%
		\begin{tabular}{lrrrrrrr}
			\toprule
			Run & Params & Tok. & Batch & Steps & Best step & Best val. loss & Best val. top-5 acc. \\
			\midrule
			\texttt{tok\_14\_batch\_16} & 3.70M & 14 & 16 & 7287 & 4000 & \textbf{1.4456} & \textbf{0.8619} \\
			\texttt{tok\_14} & 3.70M & 14 & 8 & 14574 & 6000 & 1.4471 & 0.8614 \\
			\texttt{constant\_1e-4} & 3.67M & 11 & 8 & 12705 & 12000 & 1.5759 & 0.8365 \\
			\texttt{xlstm-small-batch8-tok11} & 1.63M & 11 & 8 & 12705 & 12000 & 1.5698 & 0.8388 \\
			\texttt{xlstm-base-batch8-tok11} & 3.67M & 11 & 8 & 12705 & 10000 & 1.5811 & 0.8367 \\
			\texttt{xlstm-medium-batch8-tok14} & 9.28M & 14 & 8 & 72870 & 9000 & 1.5173 & 0.8542 \\
			\bottomrule
		\end{tabular}
	}
	\caption{Representative validation results from local experiment summaries. Lower validation loss is better. The medium row reports the early best checkpoint, not the final checkpoint.}
	\label{tab:result-checkpoints}
\end{table}

The validation ranking supports two main observations. First, tokenizer 14 is a
better downstream representation for the current xLSTM setup than tokenizer 11,
even though tokenizer 11 was the best tokenizer-only round-trip candidate. This
confirms that compactness and reconstruction metrics are useful diagnostics but
do not fully predict language-model training behavior. Second, simply increasing
model size is not enough. The medium tokenizer-14 run reaches its best validation
loss early and then overfits, so its final checkpoint is worse than its selected
checkpoint and worse than the base tokenizer-14 candidates.

\subsection{Compute Summary}

Table~\ref{tab:compute-summary} reports the logged wall-clock runtime recorded
in W\&B, together with GPU metadata from each run's \verb|wandb-metadata.json|
file. These values are approximate compute accounting numbers for the practical
work: they combine runs made during different stages of the search and should
not be read as controlled hardware benchmarks.

\begin{table}[htbp]
	\centering
	\small
	\resizebox{\textwidth}{!}{%
		\begin{tabular}{lrrlll}
			\toprule
			Experiment stage & Runs & Logged hours & GPU used & VRAM & Purpose \\
			\midrule
			Warm-up scale/batch checks & 4 & 6.78 & RTX 3050 / RTX 3060 & 6 / 12 GiB & Verify training behavior and basic scale choices \\
			Learning-rate sweep & 10 & 15.86 & RTX 3060 & 12 GiB & Compare scheduler and learning-rate settings \\
			Tokenizer sweep & 8 & 13.85 & RTX 3060 & 12 GiB & Compare REMI-style tokenizer conditions \\
			Refined tokenizer--batch sweep & 6 & 11.53 & RTX 3060 & 12 GiB & Recheck tokenizer 10/14 with batch sizes 4, 8, and 16 \\
			Medium-model follow-up & 1 & 17.76 & RTX 3060 & 12 GiB & Test larger xLSTM capacity with longer training \\
			\midrule
			Total summarized xLSTM training & 29 & 65.78 & Mixed RTX 3050/3060 runs & 6--12 GiB & Validation-stage experimental budget \\
			\bottomrule
		\end{tabular}
	}
	\caption{Approximate compute used by completed xLSTM experiments, aggregated from W\&B runtime summaries and run metadata.}
	\label{tab:compute-summary}
\end{table}

The compute distribution also explains the current confidence level of the
results. The majority of the budget was spent on validation sweeps rather than
on one final long training run. This is appropriate for the current practical
phase because the main uncertainty was representation and optimization choice.
The next evaluation step should spend compute on held-out generation from the
shortlisted checkpoints rather than on extending the same validation-only
ranking.
